\section{Discussion}
\label{sec:discussion}

\subsection{Does the quantum kernel help on this task?}
No, not on this data. Under the corrected protocol the ZZ quantum kernel matches a tuned
classical RBF kernel on the same PCA features. It does not beat it on any metric, circuit
architecture, qubit count or training-set size we tested, and a Random Forest on all raw
features beats both. This agrees with large systematic benchmarks of quantum models on classical
data~\cite{bowles2024}, which found that well-tuned classical baselines usually match or beat
them.

This is a negative result about \emph{this} dataset and feature map, not about quantum kernels in
general. The diagnostics show why it comes out this way:
\begin{enumerate}
\item \textbf{The problem is easy and low-dimensional.} Binary accuracy is 99.5\% with only 4 PCA
      components, and every model saturates at 6--8 components. There is little left for a more
      expressive kernel to find.
\item \textbf{The quantum kernel's geometry is classically reachable.} With tuned $\alpha$, the
      geometric difference to the RBF kernel is 3--14, well below $\sqrt{N}=20$. By
      Huang et al.~\cite{huang2021} a classical kernel can then reproduce the quantum kernel's
      predictions, so no large advantage is possible, whatever the labels.
\item \textbf{The useful regime of the ZZ map is the classical-looking one.} The tuning selects
      small $\alpha$ for the 16-qubit, 2-repetition circuit (0.05--0.4). There, the phases are small and the kernel behaves
      like a smooth, RBF-like similarity. The large-$\alpha$ regime, where the circuit is
      conjectured to be hard to simulate, is exactly where it concentrates and loses accuracy
      (Figure~\ref{fig:bandwidth}).
\item \textbf{Entanglement is not used.} The separable $R_x$ kernel, which is a classical closed-form
      expression, performs the same as the entangling ZZ kernel, and more repetitions or denser
      entanglement do not change accuracy (Table~\ref{tab:arch}).
\item \textbf{Flow statistics have no quantum structure.} Known quantum advantages for kernels use
      data whose labels are generated by a quantum process or encode a problem that is hard
      classically (e.g.\ discrete logarithms). CICFlowMeter statistics are classical summaries of
      packet counts and timings, and the ZZ feature map was not designed around them.
\end{enumerate}

\subsection{What went wrong in the original evaluation}
The original advantages were produced by the evaluation, not by the model. Ranked by effect:
\begin{enumerate}
\item \textbf{Duplicate leakage} (up to 49\% of test flows with a training twin). It inflated the
      QSVM more than the RBF-SVM: a nearly diagonal kernel memorises twins perfectly.
\item \textbf{Linear encoding of heavy-tailed features.} It hurt the RBF baseline much more than
      the QSVM (3 points in binary mode), creating most of the binary gap.
\item \textbf{Unequal tuning.} The RBF-SVM was run with default $C$ and $\gamma$. Comparing a
      configured model against a default one measures the configuration, not the model.
\item \textbf{Single small test sets.} The binary ZZ result rested on 165 test flows, where
      ``98.79\% vs.\ 96.97\%'' is 163 vs.\ 160 correct. No repetition or significance test was
      reported.
\end{enumerate}
These are general lessons for quantum-ML benchmarks on network data: deduplicate before
splitting, fit every transform on training data only, give the classical baseline the same
tuning budget, report several seeds with a paired test, and include a strong non-kernel
baseline.

\subsection{Threats to validity}
\begin{itemize}
\item \textbf{Simulation only.} All kernels are exact or binomially sampled statevector
      simulations. Hardware adds gate errors, decoherence and readout noise. The shot
      experiments model only statistical noise, and even that costs 3 points in the 16-qubit
      multiclass setting (Section~\ref{sec:res-shots}).
\item \textbf{One capture, contiguous slices.} Each class in \code{output.csv} is a contiguous slice of
      one capture day, in some cases spanning about a second. Random train/test splits within such
      a slice are easier than training on one day and testing on another. Cross-day or cross-dataset
      evaluation would be the stronger test, and it would make all models' accuracy lower.
\item \textbf{Balanced sampling.} Real traffic is heavily imbalanced, so the false-alarm rate would
      matter much more than accuracy. The FAR is reported, but not at a realistic base rate.
\item \textbf{Statistical power.} Each seed tests 400 flows, so per-seed McNemar tests only detect
      differences of roughly 2 points or more. Pooled over 5 seeds, the discordant counts are nearly
      balanced (2 vs.\ 1 binary, 19 vs.\ 23 multiclass), so a large hidden advantage is very
      unlikely, but differences below about 1 point cannot be ruled out.
\item \textbf{Shared preprocessing inside CV.} PCA and scaling are fitted once on the full
      training split, and the CV folds for tuning reuse that fit. The validation folds therefore
      contribute to the unsupervised preprocessing. This affects the QSVM and the RBF-SVM equally
      and does not touch the test set.
\item \textbf{Bandwidth grid boundary for $R_x$.} The tuned $\alpha$ of the $R_x$ kernel was 1.0, the
      top of the default grid, in all 10 headline runs. Rerunning the 16-qubit $R_x$ headline with
      $\alpha\in\{\dots,1.5,2,3\}$ selects $\alpha=1.5$--3 and raises accuracy slightly (binary
      99.50\% $\to$ 99.60\%, multiclass 91.05\% $\to$ 91.65\%), still with no significant seed
      against the RBF-SVM ($p\ge0.125$). The conclusion is unchanged, but the reported $R_x$
      numbers are slightly conservative. For the ZZ map the tuned $\alpha$ lies inside the grid.
\item \textbf{Ambiguous rows.} After deduplication, 177 multiclass rows (2 binary) share their
      feature vector with a row of a different class. They are kept as genuine ambiguity and
      cause the 0.8\% residual train/test overlap in multiclass mode.
\end{itemize}

\subsection{Future work}
\begin{enumerate}
\item \textbf{Harder data:} the 12 January capture day of CIC-DDoS2019 (DNS, NTP, SSDP, SNMP and
      other near-identical reflection attacks), cross-day train/test splits, and other datasets
      such as CICIDS2017 or UNSW-NB15. \code{--csv} already accepts any CICFlowMeter file.
\item \textbf{Trainable feature maps} that maximise kernel--target alignment
      (\code{evaluation.centered\_alignment} is already implemented).
\item \textbf{Projected quantum kernels}~\cite{huang2021}, which are less prone to concentration.
\item \textbf{Noise models and mitigation:} depolarising and readout noise on the Cirq backend,
      and noise-aware kernel post-processing that treats the training and test kernels
      consistently. Clipping only the training kernel's spectrum made results worse here.
\item \textbf{Hardware runs} of a subset of kernel entries using \code{build\_kernel\_circuit()}.
\item \textbf{Imbalanced evaluation:} train on balanced data, test at the natural class ratio, and
      report precision at a fixed false-alarm rate.
\end{enumerate}

\subsection{Conclusion}
The project now has a correct, fast and reproducible quantum-kernel pipeline. Kernels are
verified against Cirq, 16-qubit kernels take seconds instead of days, preprocessing is
leakage-free, tuning is equal, and the evaluation uses several seeds with significance tests.
With that pipeline, the reported quantum advantage on CIC-DDoS2019 does not survive. The ZZ
quantum kernel SVM is as good as, but not better than, a tuned classical RBF-SVM, and both are
slightly behind a Random Forest on all features. The tools built here (kernel diagnostics, the
geometric-difference test, the benchmark suite) can be used to test whether a quantum kernel helps
on harder data, and to find out quickly when it does not.
